% NAL 2338, batch 4 — Gallica views 61–80.
% Pass: Opus 5.5 (claude-opus-5-5), 2026-09-23 — first pass, unchecked against the views by a human.
% Read from the local mirror only (archives/nal-2338/f0061.jpg … f0080.jpg):
% 1000-px thumbnails first, then four full-width bands per written view at
% 2400 px, then tight crops for doubtful words, the marginal addition of
% v. 74, the root sign of v. 80 and the figure of v. 71. Each view was
% drafted to the scratchpad (nal-2338/drafts-b4/) as it was read, and this
% file assembled from the drafts. No edition of Roberval, Torricelli or
% Mersenne and no other transcription was opened; the verses quoted in the
% letter are identified in notes from memory, as the pass's own
% rapprochements, and nothing was taken from an edition.
%
% Dating. The catalogue gives none: \dating{} is left empty, and no date is
% inferred for the volume. The leaves carry no date of writing in these
% views; the years named in the text (1637, 1643, 1644, 1645, and « circa
% finem anni 1644 ») are content and are transcribed where they stand.
%
% What the twenty views hold. Greyscale and colour scans, portrait, about
% 4316 × 5740, one page per view, the edge of the next leaf and the binding
% at one side. One piece only: the continuation of the Latin letter that
% begins at v. 57 (batch 3) under the heading « Epistola Aegidii Personerii
% de Roberval ad Euangelistam Torricellium ». In these views the letter is
% addressed to Torricelli by name (« Crede mihi clarissime Toricelli »,
% v. 77) and the writer speaks as the inventor of the lines Torricelli calls
% « Roberuallianas » (v. 70) and of the trochoid; none of these views is
% signed. The writer named is the heading's, on v. 57; that the text is his
% is what the heading says, not what these views show. The hand is not his
% necessarily: the pages are a fair copy (see Hands).
%
%   v. 61–62   Numbers: sums of powers of the integers « ordine naturali,
%              atque indefinitè sumptorum » to the greatest taken as many
%              times, 1/3, 1/4, 1/5 …, and of square and cube roots, 2 ad 3,
%              2 ad 5, 3 ad 4 …; the proof « per duplicem positionem more
%              veterum »; sums of powers of the first 100000000 numbers.
%   v. 62–64   Centres of gravity: the writer's method against the one
%              Torricelli sent to Mersenne « circa finem anni 1644 »; the
%              charge of plagiarism answered by dates (the writer's method
%              « jam ab anno 1637 »); the fable of Aesop's sculptor
%              (Jupiter, Neptune, Mercury: the plane, the solid, the centre).
%   v. 65–66   Mersenne as go-between: the solid of the trochoid about its
%              axis, Torricelli's ratio 11 ad 18 found smaller than the true
%              one, 5 ad 8 about the base; a plea that no letter be trusted
%              unless signed by the writer's own hand; Mersenne and the
%              quadrature of the hyperbola (« Cleopatræ … vnionem »).
%   v. 67–69   The counter-charges: the trochoid published by Torricelli;
%              parabolas equal in length to spirals (the writer's, published
%              by Mersenne), conical spirals; the spiral Torricelli claims
%              equal to a straight line, and the loxodromes; Apollonius I, 20
%              and 21; the new quadratrices sent « vix … ante biennium ».
%   v. 70–71   The construction of the « primaria nostra quadratrix » AML
%              from a trilinear ABC and its tangents DF, EG, CK; v. 71 is the
%              figure alone, on a slip mounted on a blank leaf.
%   v. 72      Blank verso of the mounting leaf. Not transcribed.
%   v. 73–75   The quadrilinear ABCN double the trilinear ABC « per
%              infinita »; solids of revolution triple and double; the
%              secondary quadratrix BTS; other quadratrices; the quadrature of
%              any parabola, and of its solid, deduced from them.
%   v. 75–79   The Academy's « proceres » read the letter sent to
%              Torricelli; Torricelli's own words quoted and answered (the
%              quadratures of the parabolas are Archimedes's and Fermat's;
%              « Florem mittebamus, non arborem »); Archimedes defended
%              (projectiles; 11 ad 14 against 11 ad 18).
%   v. 79–80   The tone of the letter excused; then a list of what the
%              writer « ex nobis, tùm ex nostris » holds: equations, loci
%              plani, solidi, ad superficiem, « loca solida ad tres et
%              quatuor lineas », isoperimetric cylinders and cones, the cone
%              of greatest surface inscribed in a sphere of diameter 32
%              (axis 23 − √17), portions of the sphere, cylindrical surfaces.
%
% First and last views. v. 61 continues the letter begun at v. 57: v. 60
% ends « vniuersa- » with the catchword « lius » under the line, and v. 61
% opens « lius multò quàm … ». Batch 3 set the word « vniuersalius » on v. 60;
% here the page's own « lius » is transcribed, so the coordinator may drop
% one of the two. v. 80 runs on: it ends « Circa Arithmeticam, Musicam,
% Opticam, Astrono- » with « miam » under the line (end of the word, and
% probably the catchword of v. 81). The letter is not finished in this
% batch; no signature, date or address in views 61–80.
%
% Hands. One hand throughout, the hand B of batch 3 (v. 57–60): a regular,
% posed italic with accents on vowels (à, â, è, ò, ù), a b-shaped 6 in dates
% (« ıb44 », v. 64), catchwords at the foot of rectos and versos (« tissima: »
% v. 62, « gis » v. 64, « miam » v. 80; « lius » v. 60). Quotations from
% Torricelli's letters, and a few other words (« Vtinam non misissem »,
% v. 69), are written in a larger, rounder script: the foot-note of v. 57
% (batch 3) asks that what is in large characters be set in italics, so the
% copy appears to be prepared for print — an inference. The large script is
% described in notes, not marked in the text. Points of figures are in posed
% capitals. One marginal addition on v. 74, in a quicker and thinner pen
% (perhaps another hand), is set as \marginal{}. Very few corrections: a
% word drowned under an ink blot (v. 67), letters retouched in blacker ink.
%
% Spelling. The copy's own: u inside words and v at their head (« inuenire »,
% « vt », « vnicam »), ij, j in « ejusdem », « jam », « Toricellius » with one
% r, « fermat » without capital, « assymptotis », « hypothenusæ »,
% « sperrendum » (v. 62), « quadraticum » (v. 74), « trilineam » for
% « trilineum » (v. 70, 73), « spatiam » (v. 75), « occurátque » (v. 74). Nasal
% tildes (sumptũ, annũ, nequaquã, quã) and accents are kept as on the
% page. A long s is transcribed as s. « et cæt. » as written. Latin not
% normalised.
%
% Notation. Fractions 1/3, 1/4, 1/5 set $\frac13$ etc. (the page writes them
% as fractions). The root sign of « 23 − √17 » (v. 80) is a small r-shaped
% stroke, set \sqrt. No other sign of operation; ratios « vt 11 ad 18 » in
% words.
%
% Figures. One figure only in the batch, on the slip of v. 71, described in
% a note with its lettering. v. 70 and v. 73–75 refer to it.
%
% Quoted verses. v. 67, v. 78 set lines apart from the prose: « Aspice num
% mage sit nostrum penetrabile telum » and others (Virgil, Horace); the
% identifications are the pass's own and are given in notes as such.
%
% Numbers on the leaves.
% - One ink series at the top right of the rectos, continuing batch 3
%   (… 28 at v. 57, 29 at v. 59): 30 (v. 61), 31 (v. 63), 32 (v. 65),
%   33 (v. 67), 34 (v. 69), « 34 bis » (the figure slip, v. 71), 35 (v. 73),
%   36 (v. 75, with « 37 » of the next leaf showing above its top edge),
%   37 (v. 77), 38 (v. 79). It runs across leaves and slips, and its « bis »
%   for an inserted slip is a counter's, not the writer's: it behaves as the
%   library's foliation. It is ink, so, as for the other volumes, no
%   \folio{} is written until a human has looked. The mounting leaf of the
%   slip (v. 71–72) carries no number of its own. No verso carries a number;
%   the recto's figure shows reversed through each verso.
% - Numbers in the text are content: years 1637, 1643, 1644, 1645; ratios
%   1 ad 3 … 2 ad 9, 5 ad 8, 11 ad 18, 11 ad 14; 100000000; 32 and 23 − √17;
%   « 21am prop. » and « 20am prop. » of Apollonius; « prop. 2 », « prop. 3 »
%   of Archimedes's Dimensio circuli.
% - No stamp, no shelfmark in these views.
%
% The catalogue's dating is copied verbatim (empty).

% Coordinator's reconciliation (2026-09-23), after all six batches:
%   One ink series of leaf numbers runs top right of the rectos across every
%   hand and piece: 1–10 (v. 2–20), 11–20 (v. 22–40), 21–29 (v. 42–59),
%   30–38 with « 34 bis » for the figure slip (v. 61–79), 39–48 (v. 81–99),
%   49–50 (v. 101–103). It ends on « 50 » at the last written leaf, as the
%   certificate of v. 1 has it (« Volume de 50 Feuillets / 1er Octobre
%   1888 »): it is the library's count. Being ink, it gets no \folio{}, in any
%   batch. The bold series on the letter of v. 24–35 (1–11) is that letter's
%   own pagination (batch 2).
%   Seams: v. 60 ends « vniuersa- » and v. 61 opens « lius »; batch 3 had set
%   the catchword into v. 60, and now ends where the page does. v. 100 runs
%   into v. 101 (« altius ; » / « in aliis autem humilius »), the « De Vacuo
%   Narratio » (v. 97–103). The « non » sign is set « nl » in batch 1 (hand B,
%   an n with a barred ascender) and « n^o » in batch 2 (the fast cursive, an
%   n with a raised loop): two hands, two signs, left as each pass read them.

\documentclass[11pt,a4paper]{article}
% The preamble shared by every transcription of the mathematicians' manuscripts in Gallica.
%
% It rests on one idea: **uncertainty compounds, it does not resolve.** A
% transcription that smooths over a doubtful word has destroyed the only thing
% it was made for — a reader must be able to tell, at every point, what was on
% the page and what was supplied. The apparatus macros below are therefore the
% critical apparatus, not decoration.
%
% The same commands are recognised by `scripts/render.mjs`, which produces the
% site's reading view, and by `scripts/tei.mjs`. A macro added here must be
% added there too, or rendering fails loudly — which is the wanted behaviour.
%
% Comments here are English, like the rest of the repository. The documents
% this preamble sets are French, and that is deliberate: see the skills.

\ProvidesPackage{germain}[2026/09/15 Transcriptions of mathematicians' manuscripts in Gallica]

\usepackage{amsmath,amssymb,amsthm}
\usepackage{mathrsfs}
% Kept from the parent preamble so the renderer's diagram support stays
% available; these manuscripts draw no commutative diagrams, and nothing here uses it.
\usepackage{tikz-cd}
\usepackage[english,french]{babel}
\usepackage{fontspec}

% --- Greek ------------------------------------------------------------------
% The text face stays fontspec's default, Latin Modern, which has no Greek:
% XeTeX drops every glyph it lacks with a line in the log and nothing on the
% page, and the Greek words the manuscripts quote vanished from the PDFs. So
% the Greek and Greek Extended blocks get a XeTeX character class of their own,
% and entering or leaving a run of it switches the family to CMU Serif — the
% same Computer Modern design, with polytonic Greek — and back. Only the
% family changes: a Greek word in \textit or \textbf keeps its shape and series.
% The transitions to and from every other class carry the tokens that class
% has towards an ordinary letter, so French spacing before « ; » or inside
% « » still holds after a Greek word. No grouping, so no brace or \endgroup
% can come out unbalanced; maths is left alone, since XeTeX inserts nothing
% there. Kept identical in `exercices.sty`.
\makeatletter
\newfontfamily\germainGreekfont{cmun}[
  Extension=.otf, NFSSFamily=germaingreek,
  UprightFont=*rm, ItalicFont=*ti, SlantedFont=*sl,
  BoldFont=*bx, BoldItalicFont=*bi, BoldSlantedFont=*bl]
\newXeTeXintercharclass\germain@greek
\def\germain@greekrange#1#2{%
  \@tempcnta=#1\relax
  \loop
    \XeTeXcharclass\@tempcnta=\germain@greek
    \ifnum\@tempcnta<#2\advance\@tempcnta\@ne\repeat}
\germain@greekrange{"0370}{"03FF}
\germain@greekrange{"1F00}{"1FFF}
\def\germain@greekprev{\rmdefault}
\def\germain@greekin{%
  \edef\germain@greekprev{\f@family}\fontfamily{germaingreek}\selectfont}
\def\germain@greekout{\fontfamily{\germain@greekprev}\selectfont}
\def\germain@greekjoin{%
  \XeTeXinterchartokenstate=\@ne
  \@tempcnta=\z@
  \loop
    \ifnum\@tempcnta=\germain@greek\else
      \XeTeXinterchartoks\germain@greek\@tempcnta=\expandafter{%
        \expandafter\germain@greekout\the\XeTeXinterchartoks\z@\@tempcnta}%
      \XeTeXinterchartoks\@tempcnta\germain@greek=\expandafter{%
        \the\XeTeXinterchartoks\@tempcnta\z@\germain@greekin}%
    \fi
    \ifnum\@tempcnta<4095\advance\@tempcnta\@ne\repeat}
% babel-french sets its own classes, and saves and restores the interchar
% state, each time a language is selected: join again after it does.
\addto\extrasfrench{\germain@greekjoin}
\addto\extrasenglish{\germain@greekjoin}
\AtBeginDocument{\germain@greekjoin}
\makeatother

\usepackage{geometry}
\geometry{a4paper,margin=25mm,marginparwidth=28mm}
\usepackage{marginnote}
\usepackage{xcolor}
\usepackage[normalem]{ulem}
\setlength{\emergencystretch}{3em}
\usepackage{hyperref}
\hypersetup{colorlinks=true,linkcolor=black,urlcolor=black,pdfborder={0 0 0}}

% --- Batch metadata ---------------------------------------------------------
% Taken as they stand by the rendering script, which shows them at the head of
% the reading view. They fix what the file claims to cover. The macro names are
% the parent project's, so that the scripts stay shared; \folder here means the
% volume's slug — `fr-9115` — and \pages counts Gallica views.

\newcommand{\folder}[1]{\gdef\grothFolder{#1}}
\newcommand{\batch}[1]{\gdef\grothBatch{#1}}
\newcommand{\pages}[2]{\gdef\grothFirst{#1}\gdef\grothLast{#2}}
\newcommand{\foldertitle}[1]{\gdef\grothFoldertitle{#1}}

% The holder's own shelfmark, printed as they write it: « Français 9115 ».
\newcommand{\shelfmark}[1]{\gdef\grothShelfmark{#1}}

% The dating, copied verbatim from the holder's catalogue — never inferred
% here. « XIXe siècle » is the BnF's; a pass that dates a leaf from its content
% says so in a \note{}, as its own inference.
\newcommand{\dating}[1]{\gdef\grothDating{#1}}

% The status notice, printed in the title block and repeated as a diagonal
% watermark on every page of the PDF. The manuscripts are in the public
% domain; what this line declares is not a right but a quality — an
% unverified first pass by a machine. The skills write the line; a file
% without it gets no watermark, so leaving it out is visible in the output.
\newcommand{\watermark}[1]{\gdef\grothWatermark{#1}}

\gdef\grothFolder{?}\gdef\grothBatch{}\gdef\grothFirst{?}\gdef\grothLast{?}%
\gdef\grothFoldertitle{}\gdef\grothShelfmark{}\gdef\grothDating{}\gdef\grothWatermark{}

% --- The résumé -------------------------------------------------------------
\newenvironment{resume}
  {\small\setlength{\parskip}{0.45em}}
  {\par\medskip\hrule\medskip}

% --- Keywords ---------------------------------------------------------------
% In English, closing the modernised reading's résumé: the modern vocabulary
% under which to search for what the leaves construct. The manifest extracts
% them to tag the volume — this line is the single source of the tags.
\newcommand{\keywords}[1]{\par\smallskip\noindent
  {\footnotesize\textsc{Keywords} --- \textit{#1}}\par}

% --- The critical apparatus -------------------------------------------------

% The Gallica view number, in the margin. This is the anchor that lets a
% disagreement over a formula be settled by looking at *one* image rather than
% a volume of 750. The site uses it to turn the facsimile as the transcription
% is scrolled. A view is one scanned image: usually one side of a leaf.
\newcommand{\page}[1]{%
  \marginnote{\footnotesize\color{gray!70}\textsf{v.\,#1}}%
  \label{page:#1}\ignorespaces}

% The BnF's pencilled foliation, where the leaf carries it and a pass can read
% it: « 198r ». This is what the literature cites — « ff. 198r–208v » — and
% the one line that turns a citation into a view. Recorded, never inferred:
% a folio a pass cannot read is not written.
\newcommand{\folio}[1]{%
  \marginnote{\footnotesize\color{gray!50}\textsf{f.\,#1}}\ignorespaces}

% The modernised reading's coarser counterpart to \page{}: which views a
% section is drawn from.
\newcommand{\pagerange}[2]{%
  \marginnote{\footnotesize\color{gray!70}\textsf{v.\,#1--#2}}%
  \ignorespaces}

% Illegible: never guessed.
\newcommand{\ill}{\textcolor{red!60!black}{[\,\ldots\,]}}

% A doubtful reading: the word is given, and flagged as uncertain.
\newcommand{\uncertain}[1]{\uline{#1}}

% An editorial addition — a missing letter, an implied word.
\newcommand{\add}[1]{\textcolor{blue!50!black}{[#1]}}

% Struck out by the writer. What was crossed out is part of the document.
\newcommand{\struck}[1]{\sout{#1}}

% The transcriber's note, on the state of the page or the mathematics.
\newcommand{\note}[1]{\par\smallskip{\small\color{gray!60!black}\textit{[#1]}}\par\smallskip}

% A marginal note of hers, distinct from the body of the text.
\newcommand{\marginal}[1]{\par\smallskip\noindent\hspace{1em}%
  {\small\color{gray!60!black}#1}\par\smallskip}

% --- Title ------------------------------------------------------------------
% The title block is French, like the documents themselves.

\AddToHook{shipout/background}{%
  \ifx\grothWatermark\empty\else
    \begin{tikzpicture}[remember picture, overlay]
      \node[rotate=52, scale=3.4, align=center, text=gray!18,
            font=\sffamily\bfseries]
        at (current page.center) {\grothWatermark};
    \end{tikzpicture}%
  \fi}

\AtBeginDocument{%
  \begin{center}
    {\large\bfseries Manuscrits de mathématiciens (Gallica) ---
      \ifx\grothShelfmark\empty\grothFolder\else\grothShelfmark\fi}\\[2pt]
    {\itshape\grothFoldertitle}\\[4pt]
    {\small vues \grothFirst--\grothLast
      \ifx\grothBatch\empty\else\ (lot \grothBatch)\fi}\\[2pt]
    \ifx\grothDating\empty\else
      {\small\color{gray!60!black}Datation du catalogue : \grothDating}\\[2pt]
    \fi
    \ifx\grothWatermark\empty\else
      {\small\bfseries\color{red!45!gray}\grothWatermark}\\[2pt]
    \fi
    {\footnotesize\color{gray!60!black}Transcription établie d'après les images IIIF de Gallica.
     Source gallica.bnf.fr / Bibliothèque nationale de France.\\
     Les lectures incertaines sont soulignées, les illisibles notées [\,\ldots\,],
     les ajouts entre crochets ; \textsf{v.} numérote les vues de Gallica,
     \textsf{f.} la foliotation au crayon de la BnF.}
  \end{center}
  \vspace{1em}}

\endinput


\folder{nal-2338}
\shelfmark{NAL 2338}
\batch{4}
\pages{61}{80}
\dating{}
\watermark{Lecture automatique\\non vérifiée}
\foldertitle{Papiers divers de Mersenne.}

\begin{document}

\page{61}
\note{La page commence par la fin d'un mot, « lius » : au bas de la vue 60 (lot 3), « vniuersa- » est suivi, sous la ligne, de « lius », réclame que cette page reprend ; le mot est donc « vniuersalius », et la phrase de la lettre commencée à la vue 57 continue ici. Même main qu'aux vues 57--60 : italique posée, serrée, avec les accents (à, ò, ù, è) et le 6 en forme de b dans les millésimes. En haut à droite, à l'encre, « 30 » (voir l'en-tête). Le bas de la page est blanc sous la dernière ligne, qui s'arrête sur une virgule.}

lius multò quàm ipse proponas: quippe non solùm potestates in helicibus proposuit, sed etiam potestatum radices. Exempli gratiâ. Si in helice semidiametri omnium reuolutionum ordine sumptarum, se habeant vt radices quadratæ, aut cubicæ et cæt. numerorum ordine naturali progredientium 1, 2, 3, 4, 5, 6, 7, 8, et cæt. quarum primam (quadraticam putà) reperies in prima reuolutione dimidiam partem sui circuli constituere. Cúmque ipsum arduarum (vt tunc) propositionum demonstrationes rogarem, ille in hæc verba rescripsit. Ego, inquit, vt inuenirem laboraui; labora et ipse: in hoc enim labore præcipuam voluptatis partem consistere deprehendes. Quid facerem à tanto viro incitatus? laboraui, atque in auxilium infinita nostra aduocaui; (nondum enim tunc nostra ampliùs non esse resciueram) eáque tum primùm ad numeros extendi.

Animaduerti enim, et parabolarum plana, ad sua parallelogramma; et earundem solida, ad suos cylindros; et spatia helicum, ad suos circulos feliciter comparari posse, si innotesceret in numeris ratio summæ potestatum omnium ejusdem generis, ordine, atque indefinitè sumptarum, ad earum maximam totiès sumptam; ídque in omni genere potestatum. Quod quidem non difficulter assecutus sum: illicò enim patuit summam omnium numerorum quadratorum, ordine naturali, atque indefinitè sumptorum, 1, 4, 9, 16, 25, et cæt. ad eorum maximum totiès sumptum quot sunt illi quadrati; hoc est ad cubum ejusdem radicis cum maximo illo quadrato, collatam, se habere vt 1 ad 3, siue constituere $\frac13$. Summam cuborum eodem modo sumptorum, ad eorum maximum totiès sumptũ, siue ad quadrato-quadratum ejusdem radicis cum maximo cubo, se habere vt 1 ad 4; siue constituere $\frac14$. Summam quadrato-quadratorum, eodem modo, constituere $\frac15$. Atque ita in infinitum. Ex hac propositione quæ sola sufficit, innumera deduxi corollaria, qualia sunt hæc. Summa radicum quadratarum numerorum omnium, ordine naturali, atque indefinitè sumptorum, ad earundem radicum maximam totiès sumptam, collata; putà Summa radicum quadratarum horum numerorum 1, 2, 3, 4, 5, 6, et cæt. eam habet rationem quam 2 ad 3. Summa radicum quadratarum omnium numerorum quadratorum, ordine naturali, atque indefinitè sumptorum, ad earundem radicum maximam totiès sumptam, se habet vt 2 ad 4. Summa radicum quadratarum omnium numerorum cuborũ, ad maximam totiès sumptam, vt supra, se habet vt 2 ad 5. Atque ita in infinitum, radices quadratæ numerorum quadrato-quadratorum, quadrato-cuborum, cubo-cuborum, et cæt. ad earum maximam totiès sumptam, vt supra, sic comparabuntur, vt antecedens rationis sit semper 2 exponens quadrati; consequens verò sit summa ex ipso exponente 2 et alio exponente ipsius gradus ad quem pertinent numeri quorum sumuntur radices quadratæ. Vt si sumantur radices quadratæ numerorum quadrato-quadrato-cuborum qui sunt septimi gradus cujus exponens est 7, erit consequens rationis 9, conflatum, ex 2 et 7: et ratio erit vt 2 ad 9. Similiter summa omnium radicum cubicarum omnium numerorum ordine naturali,
\note{La phrase reste en suspens au bas de la page ; elle continue à la vue 62.}

\page{62}
\note{Verso du feuillet « 30 », même main ; la phrase de la vue 61 continue. Aucun numéro sur la page. Le bas de la page est blanc.}

hoc est in primo gradu, atque indefinitè sumptorum, ad earundem radicum maximam totiès sumptam, se habet vt 3 exponens cubi, ad 4 compositum ex eodem 3 et 1 exponente primi gradus. Summa omnium radicum cubicarum omnium quadratorum, ad earundem radicum maximam totiès sumptam, vt supra, se habet vt 3 ad 5. Atque ita in infinitum, radices cubicæ omnium graduum, ad earundem maximam sumptam vt supra, comparabuntur; eritque in omnibus antecedens 3, consequens verò componetur ex eodem 3 juncto cum exponente gradus cujus radix cubica sumpta fuerit. Nec aliter radices quadrato-quadratæ omnium graduum, ad earum maximam sumptam vt dictum est, comparabuntur; eritque antecedens 4. Et sic in infinitum infinitiès, vt satis ex prædictis patet. Hæc cùm ad amplissimum virum scripsissem, dubitauit num eorum demonstrationes haberem. Itaque paucis verbis indicaui eam esse facillimam, per duplicem positionem more veterum, incipiendo ab vnitate, et procedendo ordine per omnes potestates. Quo pacto, facilè est concludere in quadratis, exempli gratiâ; Summam omnium numerorum quadratorum ordine naturali, sed finitè, sumptorum, ad eorundem maximum totiès sumptum, collatam, majorem esse quàm $\frac13$; at dempto ab eadem summa, seu ab antecedente rationis, ipsorum quadratorum maximo tantùm, remanente integro consequente, relinqui rationem minorem quàm $\frac13$. Nec ad id demonstrandum, aliò recurrendum est quàm ad genesim quadratorum, quâ fit vt quiuis numerus quadratus componatur ex proximo quadrato minore, ex duplo radicis ejusdem minoris, atque ex vnitate. Quemadmodum etiam quiuis numerus cubus componitur ex proximo cubo minore, ex triplo quadrati minoris, ex triplo radicis minoris, atque ex vnitate. Qui quidem cubus est ipsum maximum quadratũ totiès sumptum quot sunt numeri quadrati ab vnitate incipientes: Atque ita de singulis potestatibus, secundùm vniuscujusque genesim. Corollaria, quomodò ab ijs deducantur, aliàs, si ita expediat, explicabimus.

Neque etiam fortassis sperrendum videbitur corollarium aliud quod ex tali numerorum inspectione deduxi: illud autem tale est. Propositis quotcunque numeris multitudine finitis, qui ab vnitate, secundùm naturalem numerorum seriem procedant 1, 2, 3, 4, 5, 6, 7, 8, et cæt. vsque ad 100000000 exempli gratiâ. Exhibere summam quadratorum, aut cuborum, aut quadrato quadratorum, aut cubo-quadratorum, aut cubo-cuborum, et cæt. omnium talium numerorum. Quæ sanè regula, pro quadratis, et cubis, reperitur specialis apud authores: at pro omnibus potestatibus, nullam apud illos reperimus vniuersalem. Hæc ergo fuit nostra pro parabolarum planis ac solidis, simúlque pro planis helicum, methodus. Post hæc proposuit vir amplissimus (quod et ipse jamdiu in omnibus figuris vniuersaliter quærebam) prædictarum figurarum centra grauitatis inuenire. Ac ille quidem ad Analysim recurrit, nos ad nostra infinita. Vnde methodus illius, vt plerisque inuentis analyticis accidit, abstrusissima est, subtilissima, atque elegan-
\note{Au pied de la page, seul contre la marge droite, « tissima: », réclame de la page suivante, qui commence par ce même mot (vue 63). Plus haut, « sperrendum » est ainsi sur la page (pour « spernendum »).}

\page{63}
\note{Même main. En haut à droite, à l'encre, « 31 » (voir l'en-tête). Le premier mot, « tissima: », achève le mot coupé au bas de la vue 62 (réclame). Le bas de la page est blanc sous la dernière ligne, qui s'arrête sur un point-virgule.}

tissima: nostra aliquot mensibus posterior, simplicior euasit, et vniuersalior: quò fit vt cæteris collata, magis nobis arrideat. Vt tamen alicui possit esse vniuersalis, debet is omnibus numeris absolutus esse Geometra, qualis huc vsque nullus apparuit. Quoniam verò hoc nostræ hujusce dissertationis præcipuum caput est: ac vos non obliquè aut occultè, sed directè et apertè innuistis methodum nostram, quam tamen huc vsque nondùm vidistis, illius quam circa finem anni 1644 ad R. P. Mersennum à vobis missam legimus, esse inuersam; ac proinde nostram à vestra fuisse desumptam: quo posito tanquam vero, adeò indignamini, vt tres maximas epistolas ad amplissimos celeberrimósque viros, adjectis etiam ad id, magnis appendicibus, grauissimis querelis impleueritis; quò nos nihil tale meritos, acerbissimâ plagiarij contumeliâ afficeretis. Idcirco, et locus, et res postulat vt tam atrocem injuriam; quandoquidem et licet, et facilè possumus, à nobis propellamus. Ad hoc autem satis supérque futurum speraui, si nostram illam methodum ad vos cum demonstratione mitterem; non quidem suis omnibus numeris absolutam, nimis enim longa est, sed sic digestam, vt à vobis, alijsque non vulgaribus Geometris nullo negotio intelligatur; præcipuè ab ijs qui indiuisibilia non oderint: alios enim nihil moror, et Geometrarum nomine indignos puto, qui viâ apertâ, tutâ, atque facili, relictâ, longos ac difficiles anfractus sequi malint. Hoc pacto, cùm illa nostra à vestra planè diuersa sit, ac diuersis omninò fundamentis innitatur, non erit ampliùs quòd vobis ereptam conqueri jure possitis. Eam ergo seorsim cum suis figuris conscripsimus, ne hujus epistolæ lectionem interturbaret.

Facile autem erit animaduertere methodum illam eo modo quo proposita est, vniuersalem quidem esse absoluto Geometræ, attamen eandem à priori rarò procedere (vniuersalem autem à priori inuenire, hoc est ex sola figuræ aut lineæ definitione, nullâ ejus cum aliâ quauis figurâ, aut Lineâ, comparatione factâ, vix sperandum puto: quæ tamen si haberetur; et circuli, et hyperbolæ, aliarúmque numero infinitarum figurarum quadratura simul haberetur) siquidem illa in figuris, vix solâ plani cum plano, aut solâ solidi cum solido comparatione contenta, vtramque simul, et plani, et solidi, aut etiam altioris speciei comparationem persæpe requirit. Imò, illâ methodo, solidorum centra vix directè, sed plerumque indirectè tantùm, putà mediante aliquo plano congruo, deteguntur. Sed nec illa linearum centris inseruit, nisi ipsæ lineæ, earúmque proprietates quædam ex præcipuis ac specificis examinari geometricè possint. Quæ omnia ex adjectis exemplis post ipsam methodum, seorsim videre licet. De methodo domini de fermat, nisi eam adhuc videris, hoc scies, ipsam trianguli, atque planorum parabolicorum omnium et solidorum ab ijs ortorum centra à priori elegantissimè ostendere. Verùm eandem aliarum figurarum centris accommodare, hic labor; cùm ne quidem à posteriori, reliquis figuris huc vsque inseruierit;
\note{Le « F » de « Facile » est une grande capitale ornée, précédée d'un long trait de plume dans la marge. « domini de fermat » : « fermat » est écrit sans capitale.}

\page{64}
\note{Verso du feuillet « 31 », même main ; la phrase de la vue 63 continue. Aucun numéro sur la page. Le bas de la page est blanc.}

quamquam forsan, quò minus id fieri possit, nihil repugnet. Jam quòd ad tempus attinet, meministi opinor, vir clarissime methodum vestram non ante annũ 1644 Parisios missam fuisse, atque eandem tunc admodùm recens inuentam: siquidem, vt ex vestris literis patet, vobis eâ adjutis, solidi Trochoidis circa basim mensura paulò ante demùm patuerat, quam sub finem anni 1643 nondum habebatis: hæc enim sunt vestra verba in prima vestrarum ad me epistola Quoad solida; nihil habeo. Ego verò meâ methodo vsus sum jam ab anno 1637, atque illius ope, et planorum parabolicorum omnium, et solidorum centra jam tum inueneram; quorum centrorum, quæ ad dimidios fusos parabolicos pertinent, enuntiaui eâ epistolâ quam ad R. P. Mersennum de vestris inuentis scripsi anno 1643. quo primùm anno de Toricellio Parisijs auditum est. Hæc inquam enuntiaui anno plusquàm integro priusquàm vestra illa methodus appareret; quæ vestris forsan, et nostris, vnà cum aliorum inuentis (ingeniosè procul dubio) collatis, tandem apparuit. Sed finge, id quod non est, ipsam vestrã ante annum 1644 fuisse inuentam. finge etiam, id quod multò magis non est, ipsam cum nostra prorsus conuenire, ac planè eandem esse. Quid tum? An nos nostram statim vt minimè nostram repudiabimus, qui eâ septennio integro ante prædictum illum annum 1644 tanquam nostrâ, imò verè nostrâ, nemine reclamante, vsi fuerimus? num potiùs præscriptionis jure nos tutabimur? Et quibuscunque intercedentibus, nostram vt nostram lege asseremus? cùm in talium rerum possessione, vel vnius diei præscriptionem valere, nemo inficiari possit? Multò ergo potiori jure nunc, quandoquidem nostra et tempore longè prior est, et penitùs diuersa, intercessoribus valere jussis, et nostra tota manebit, qualiscunque tandem illa sit; et nostram vbique asserere, et fructibus ab ea productis tanquam nostris vti vbique licebit. Sed neque argumenta quæ produxisti, ejus ponderis esse videntur, vt illa quemquam ex ijs qui nos vel mediocriter norunt, in tam sinistram de nobis opinionem pertraherent. Primùm enim, dum ais me nunquam ne verbum quidem fecisse de centro grauitatis Trochoidis; cùm intereà tantoperè, et quidem meritò, gloriarer de omnibus alijs, quadraturâ, (comparationem cum circulo dicere voluisti) tangentibus, solidis, et cæt. Nec verissimile esse, cùm reliqua omnia proponerem, de vnico centro grauitatis siluisse, si illud tantùm sperauissem: quod quidem problema, tuo judicio, nulli reliquorum posthabendum videtur. Dum hæc ais, inquam, vir clarissime ex tuo genio loqueris: nos, dum scripsimus, ex nostro etiam genio scripsimus. Tu, cùm magnifaceres centra, quia ex ijs solida deducere posse confidebas; solida autem præcipuè intendebas; ideò centrorum inuentionem magnificè extulisti, nec cæteris posthabendam, imò præhabendam judicasti. Ego contrà, quia sine centris, solida et quæsiui, et viâ Geometricâ inueni; datis autem solidis, statim, et absque labore centra sequebantur; ideò centra ne respexi quidem, neque ad ea vnquam animum applicui; certus omninò, ex præmissa nostra methodo, dato plano quod dudum habebam, sola solida mihi quærenda superesse; centra autem simul cum plano et solidis haberi. Quòd si Apologo vti liceat, Ego sim Æsopi illius Phry-
\note{Les millésimes 1644, 1643, 1637 sont écrits avec le 6 en forme de b (« ıb44 »), comme aux vues 57--60. « Quoad solida; nihil habeo » est écrit d'une écriture un peu plus grande : c'est la citation de la lettre du destinataire. Dans « posthabendam », les lettres « am » sont repassées d'une encre plus noire. Au pied de la page, seul contre la marge droite, « gis », réclame de la page suivante, qui commence par « gis statuarius » (vue 65).}

\page{65}
\note{Même main. En haut à droite, à l'encre, « 32 » (voir l'en-tête). « gis » achève « Phry- » du bas de la vue 64. Le bas de la page est blanc ; la phrase s'arrête sans ponctuation et continue à la vue 66.}

gis statuarius: plani Trochoidis mensura, esto mihi summi Iouis statua: mensura solidi, statua Neptuni: centrum autem, esto statua Mercurij. Iam, adsit nobis è cœlo sub formâ hominis ignoti Mercurius ipse Iouis et Maiæ filius: interrogétque quanti statua Iouis. Indicabo sanè ego, alicujus pretij. Interroget deinde de statua Neptuni. Ego et ipsam alicujus pretij indicabo. Tandem interroget de sua ipsius Mercurij statua. Quid ego? Quid autem aliud nisi hoc? amice, si priores illas duas emeris, tum tertiam hanc auctarium tibi dabo. Itaque, vir clarissime, quæ tibi Iouis aut Neptuni statuâ meritò fuit, illa nobis Mercurij tantùm statua extitit. Ignosce si placet, stilo; hoc vsi sumus vt mentem nostram aperiremus. De R. P. Mersenno, quid scripserit in ea epistola cujus verba totiès repetita contra me adducis, nescio: quid autem illi dixerim ego planè memini, nec ipse omninò oblitus est: nec etiam illa quæ dixi malè congruunt cum ijs quæ sæpius pro te citasti. Sed rursus nos ex mente nostra locuti sumus; ille, vt intellexit, sic scripsit; vos ex mente vestra interpretati estis: ac illa vestra interpretatio à nostra mente alienissima est. Omnibus tamen attentè consideratis, pace tuâ dixerim vir clarissime, censui præcipuam malæ interpretationis culpam in vos recidere: neque enim verba illius, quæ ipse adducis, à nostro sensu adeò aliena fuerunt, quin ab ijs verum illum nostrum sensum facilè perspexisses, si æqui interpretis personam tibi assumere voluisses. Scripseras ad ipsũ te vtrumque Trochoidis solidum beneficio centrorum priùs inuentorum, detexisse: ac illud quidem quod circa basim, vt se habet reuerà, enũtiaueras vt 5 ad 8 quod ille cùm verum sciret (jamdudum enim ego illi tale indicaueram) non ægrè persuasus est, et alterum quoque circa axem tale esse quale affirmabas vt 11 ad 18. Lætus itaque statim ille mihi per literas significauit habere se quod mecum communicare vellet. Adiui, Epistolam tuam legi, ac circa illud postremum solidum tantùm, quod circa axem, immoratus sum; quippe quod nondum habebam, nisi in terminis vero admodùm proximis, extra quos excurrebat ratio illa à vobis assignata 11 ad 18. Hinc ergo, quia de nostris terminis nullum nobis supererat dubium, illicò animaduertimus rationem illam vestram 11 ad 18 verâ esse minorem. Cùm igitur super hâc re cogitabundus hærerem: tum R. P. ad me prior. Quid ergo, inquit, dices de clarissimo Toricellio? nonne insignium adeò theorematum cognitionem ipsi te debere fateberis? faterer, respondi, si vera essent; at talia non esse certus sum: miror sanè quòd vir talis falsum pro vero nobis velit obtrudere: nec aliud suspicari possum, nisi quòd ille mechanicâ quâdam ratione, per approximationem, hujusmodi rationem à vero non admodum longè aberrantem inuenerit, existimaueritque veram rationem non posse detegi; ac proinde suam haud veram esse, à nemine posse demonstrari. Hæc inquam ego tum, oratione fateor, planè Scyticâ; quam ille suâ ad vos epistolâ leniuit, pro suo genio, qui omninò mitis est, vt ex stilo ejus satis perspicere potuistis. Iam, cùm dixi, faterer me debere, si vera essent; planum est me non intellexisse de solido

\page{66}
\note{Verso du feuillet « 32 », même main ; la phrase de la vue 65 continue. Aucun numéro sur la page (le « 32 » du recto transparaît, inversé, en haut à gauche). Le bas de la page est blanc.}

circa basim quod jamdiu ante vos habebam, et habere me ad vos scripseram: neque de centro Trochoidis, quod dato tali solido, vnà cum plano, latere non poterat. Intellexi ergo de solido circa axem, ac de centro hemitrochoidis quod ab eo dependet, quæ etiamsi breui habiturum me confidebam, tamen jure præscriptionis, vestra fuissent, si vestra illa enuntiatio cum vero congruisset. Hinc sanè nemo non videt minimè difficile fuisse, ex verbis epistolæ R\textsuperscript{i} P\textsuperscript{is} quæ vos totiès citauistis, verum sensum qualem jam attulimus, elicere. Sed nescio quo fato aliter accidit; vnde lis hæc, pro re nullius ferè momenti, putà pro nugis nostris, vt ipse sæpe loqueris, inter nos suscepta est. Itaque, ne quid in posterum simile accidat, si tale commercium inter nos continuetur, oro vos vbicunque agetur de propositione mathematica cujus discussio ad me pertinebit, ne cujuscunque literis fidem habeatis, nisi manu meâ illæ obsignata sint: Sic enim fiet vt ego mea tantùm, non etiam aliorum scripta, ex meo sensu interpretari tenear. Nam, pace amicorum hoc dictum esto, hac in materia, soli mihi fidere assueui, jamdudum expertus, interpretes plerosque, vel dum amicis blandiri appetunt; vel dum rem non satis intelligunt, omnia literis obscurare ac prorsùs deformare. vnde qui tales literas accipiunt, illi, dum vel placitis laudibus ac blanditijs auidè sese ingurgitant, vel quod obscurum est ad placitum sibi sensum detorquent, fit necessariò vt et scribentis, et primi authoris verum sensum longè relinquant. Ac hujusmodi quidem allucinationis exemplum afferam ex tuis ipsius literis, ex proprio tuo sensu, sine interprete ad R. P. Mersennum scriptis, in quibus hæc habes. Tibi verò vir clarissime corollariolum mitto ex ipsis hyperbolis deductum Quadratura quædam est, quarum centenas, imò infinitas poteram mittere, nisi vidissem satis supérque esse vnam, vt statim omnes emergant. deinde in ijs quas ad nos scribis, quas ipse R. P. etiam ante nos legerat, hæc habes. Si vnius hyperbolæ primariæ quadratura tam diu quæsita est, nos pro vna infinitas damus. Ex quibus verbis statim existimauit R. P. primariæ hyperboles quadraturam à te inuentam fuisse. Itaque cùm aliquo post tempore, de ipsis quadraturis cum eo colloquerer, dicerémque non difficulter illas assecutum esse me. Habes ergo tandem, inquit ille, hyperbolæ conicæ quadraturam. Nequaquam, respondi, neque enim legitima hæc et nothæ illæ ijsdem Legibus addictæ sunt. Me misellum, inquit, quantâ spe decido! qui vbi Cleopatræ aut etiam majoris pretij vnionem speraui, ibi vitreas tantùm ampullas reperio. Sed de hoc ipse forsan rescribet: ego verò ideò scripsi, vt tali exemplo monerem, hac in materia non esse tutum interprete vti; cùm etiam absque hoc, tantæ eueniant allucinationes. His ergo nostris rationibus, acerbissimæ vestræ accusationis argumentis luculenter respondisse, atque cumulatè satisfecisse speramus. Nunc verò.
\note{Les deux citations des lettres du destinataire, « Tibi verò vir clarissime \ldots{} omnes emergant » et « Si vnius hyperbolæ \ldots{} infinitas damus », sont écrites d'une écriture plus grande. « R\textsuperscript{i} P\textsuperscript{is} » : « Reuerendi Patris » abrégé. La page s'achève sur « Nunc verò. » ; la suite est à la vue 67.}

\page{67}
\note{Même main. En haut à droite, à l'encre, « 33 » (voir l'en-tête). La page continue le « Nunc verò. » de la vue 66. Le bas de la page est blanc sous la dernière ligne, qui s'arrête sans ponctuation.}

Aspice num mage sit nostrum penetrabile telum.
\note{Ligne isolée, centrée : un vers cité (Virgile, \emph{Énéide}, X, 481 ; l'identification est celle de la transcription, la page ne nomme pas l'auteur).}

Videamus inquam nunc, num sit quod de vobis multò potiori jure queri possim. Ac primùm. Nonne vos Trochoidem nostram, postquàm et à R. P. Mersenno et à nobis moniti estis, jam à multis annis eam nostram esse, eámque breui à nobis in lucem emittendam; postquàm vestris ad ipsum R. P. et ad me literis polliciti estis vos talem messem nobis relicturos intactam; tamen omni jure, ac vestrâ etiam fide violatis, tanquam vestram non literis modò manu scriptis (quamquàm neque hoc ferendum fuerit) sed libello ad id prelis commisso, vulgauistis? idque interim, ac eodem prorsùs tempore quo continuis vestris literis contraria promitteretis? hæccine vestra religio? hæc consuetudo? Quòd si ego huc vsque de tali injuria, pro rei acerbitate, \struck{\ill{}} questus non sum; fateor, soli ne id facerem euicerunt communes amici: quid autem lucri feci illis obtemperando? nempe creuit vobis fiducia, quia me bardum, qui illatarum injuriarum nihil sentirem, existimauistis. Attamen si ad paucula verba quæ super hâc re ad vos scripsi, animum aduerteritis, facilè ex ijs percipietis de me dici posse
\note{À la fin de la ligne, après « acerbitate, », un mot commencé (un p ou un q initial, suivi d'un trait d'union) est noyé sous une rature d'encre noire ; « questus » est récrit en entier au début de la ligne suivante.}

Vultu simulat: premit altum corde dolorem.
\note{Ligne en retrait, à droite : un vers cité (cf. Virgile, \emph{Énéide}, I, 209, « spem vultu simulat, premit altum corde dolorem » ; rapprochement de la transcription).}

Nonne ergo ipse prior idem quod vos, sed non absque causa, clamare debui. vim patior, incredibile est quanto desiderio expectem responsum super hac re. quibus sanè verbis, ac multò etiam pluribus cùm ad R. P. Mersennum tum ad amplissimum D. de Carcauy scriptis, non obscurè significauistis vos, nisi coram vobis purgati fuerimus, in nos acerbius quidpiam omninò statuisse; vt sic et injuriâ, et multâ simul afficeremur. Sed de hoc satis: nunc ad alia capita transeamus.
\note{« vim patior \ldots{} super hac re » est écrit d'une écriture plus grande : c'est la citation des lettres du destinataire.}

Rursùs igitur, nonne primus omnium parabolas ego cum helicibus comparaui secundùm longitudinem? Nonne jam annus quintus excurrit, ex quo tale theorema vulgaui, idémque meo nomine prelis mãdauit R. P. Mersennus? nonne vos ab amicis resciuistis; ac tum demũ anno 1645 ad id animum applicuistis? habeo sanè super hâc re vestras ad vestros amicos Romanos literas vestrâ manu ac vestro idiomate scriptas. Quid tum? Jam vos palam, omnibus ferè vestris literis gloriamini, non solùm parabolam conicam cum helice Archimedea comparasse; sed et reliquas parabolas cum proprijs suis helicibus: imò et quemlibet helicis arcum vel partem, siue ex centro incipiat siue non; et siue primam reuolutionem excedat, siue non, demonstrasse cuidam lineæ parabolicæ esse æqualem. Quid hoc rei est? gloriaris de rebus nostris tanquam si tuæ illæ sint; atque id postquam nostras esse sic resciuisti, vt nisi resciuisses, ne quidem de illis, forsan, vnquam somniasses. Nec est quòd fingas existimasse te nos solam helicem Archimedeam considerasse; nimis enim frigidum fuerit figmentum, et absque vllo fundamento; cùm vna eadémque sit illius et cæterarum, demonstrationis via et methodus, quam qui inuenerit, omnia procul dubio inuenerit, si modo voluerit, nempe hæc. Quæuis parabola vnà
\note{Dans « Quid hoc rei est? », « hoc » est d'une encre plus noire, repassé.}

\page{68}
\note{Verso du feuillet « 33 », même main ; la phrase de la vue 67 continue. Aucun numéro sur la page. Le bas de la page est blanc ; la phrase s'arrête sans ponctuation et continue à la vue 69.}

cum helice sibi propriâ sic se habet, vt si portio axis parabolæ, comprehensa inter ordinatim applicatam ad axem, et tangentem à termino applicatæ ductam, æqualis esse intelligatur circumferentiæ circuli primæ reuolutionis in helice: (intellige helices planas: nos enim conicas quoque cum parabolis comparauimus) applicata autem æqualis semidiametro ejusdem circuli: tum, quæ inter verticem et applicatam interjicitur parabola æqualis sit longitudine helici primæ reuolutionis. Quòd si in eadem parabola sumatur à vertice quæuis portio; à principio autem helicis propriæ sumatur etiam portio, à cujus termino ducta recta ad helicis centrum æqualis sit rectæ à termino sumptæ portionis parabolæ, ad axem applicatæ: erunt et hæ portiones æquales. His sic à nobis inuentis, si quis quidpiam addiderit; aut si imitando, similia effecerit, habeat sanè quam ipse laudem merebitur. In helicibus conicis, existente cono recto, omnia se habent vt supra; modò tantùm loco semidiametri circuli primæ reuolutionis, qui circulus in ipso cono existit, sumatur recta à vertice coni ad circumferentiam ejusdem circuli terminata. Hîc autem, centrum helicis erit vertex coni; et quæ à centro ad puncta helicis ducuntur rectæ, erunt portiones laterum coni ejusdem. At equidem resciuisse me fateor, dices; verùm demonstrationem proprio marte adinueni. Esto: quid inde? sanè si quæstionem proposuissem tantùm, non etiam soluissem, illa tua fuisset, qui prior soluisses. Nunc quando prior solui ego, et solutam vulgaui, mea est, nec mihi, etiamsi omnes conentur, verè eripi potest. An, quæso, mea aut etiam vestra sunt parabolarum Domini de fermat quadraturæ; aut spatiórum helicum cum circulis comparationes, quas ambo proprio marte inuenimus? Quid de ipsis speretis vos, nescio sanè; ego certè, quamquam mea multò quàm vestra potior sit causa, ipsam tamen prorsùs desero. An meum est solidum vestrum hyperbolicum? an mea hyperbolarum vestrarum nouarum quadratura? minimè verò; attamen amborum ipsorum theorematum demonstrandorum vna eadémque est methodus, quam nos inuenimus, et jam pridem ad vos misimus vestro solido accommodatam, quámque ijsdem hyperbolis accommodare non admodùm difficile est. Reperi quoque in illarum singulis, ex parte vnius tantùm ex assymptotis, resecari posse spatium planum acutum et versùs acumen, infinitum, quod tamen spatio finito atque vndique clauso sit æquale. Obiter autem, vt verum fateare, nonne istis hyperbolis occasionem dedere parabolæ illæ Domini de fermat? Nonne etiam illa nostra propositio de helicibus, et parabolis longitudine æqualibus, ansam præbuit illi alteri de qua adeò magnificè gloriaris? de illo inquam helicum genere, quæ describuntur, dum recta vniformiter quidem circa manens centrum circumuoluitur, at punctum interim secundùm illam rectam fertur proportionaliter? quam quidem helicem rectæ cuidam asseris æqualem. Quæ autem sit illa recta, et quomodo ad datas se habeat, tanquam
\note{« fermat » est écrit sans capitale, les deux fois. Dans « occasionem », les lettres « ca » sont repassées d'une encre plus noire. « assymptotis » : ainsi sur la page.}

\page{69}
\note{Même main. En haut à droite, à l'encre, « 34 » (voir l'en-tête). La phrase de la vue 68 continue. Le bas de la page est blanc ; la phrase s'arrête sans ponctuation et continue à la vue 70.}

si Cereris sacrum sit, planè reticuisti: non tamen nos latet, eam æqualem esse hypothenusæ cujusdam trianguli rectanguli, cujus vnum laterum æquales sit rectæ à centro ad terminum helicis ductæ: sed enim, quis triangulum istud dabit, ex hypothesi quòd dentur positione et longitudine duæ ex ijs rectis quæ à centro ad helicem terminantur? vel contrà, quis triangulo dato, dabit helicem? vtrumque si dederis, vir clarissime. vel alterutrum tantùm, ego munus id, eo munere compensabo, quod vel ipse duplo pluris facias. Sed caue: hîc via præceps est, et lubrica; ac talis, ex qua ad paralogismum lapsus sit facillimus. Nisi tamen quod petimus datum fuerit, propositio nullius pretij remanebit. Illud etiam non videris animaduertisse, propositionem hanc non esse nouam, sed ipsam prorsùs eandem esse cum antiqua illa, quâ quæritur linea per quam pondus ad centrum terræ laberetur secundùm vniformem ad suum horizontem inclinationem: talis enim linea ad tale genus pertinet. Quàm verò minimè noua sit propositio, testabitur ipse R. P. Mersennus. Verùm, quia datâ inclinatione, hoc est, dato specie triangulo rectangulo, datoque centro helicis in centro terræ, dato insuper, vno ejusdem helicis puncto, putà in ipsius terræ superficie; non poterat geometricè, nec etiam suppositâ circuli quadraturâ, assignari aliud in ea punctum; ideò illa inculta permansit, ac ferè ex toto neglecta est. Neque rursùs, idem solum, aut primum genus est earum helicum, quæ finitæ cum sint, infinitas tamen circa punctum quoddam reuolutiones absoluunt: tales enim et longè antiquiores sunt illæ quæ in globis terrestribus, atque in mappis mundi, loxodromias, seu ventorum vias referunt: quæque præter has, illud habent peculiare, quòd ex vtraque parte finitæ sint; et tamen circa vtrumque polum infinitiès circumuoluantur. Cúmque sic imitando, res Geometricæ in infinitum plerumque abeant; quidni etiam linea recta circa manens centrum æqualiter, vel proportionaliter circumuoluetur, ac simul punctum mobile, vel æqualiter, vel inæqualiter secundùm rectam eandem, legibus quibusdam, feretur vel à centro, vel versùs centrum, ad describenda infinitiès infinita helicum genera? Ex ijs autem, genus illud nouimus, cujus helices hyperbolis conicis demonstrantur æquales. Quidni rursùs licebit, pro infinitis hyperbolis effingendis, imitari 21\textsuperscript{am} prop\textsuperscript{em} lib\textsuperscript{i} primi con\textsuperscript{rum} Apollonij, sicuti pro infinitis parabolis 20\textsuperscript{am} prop\textsuperscript{em} imitatus est D. de fermat? Verùm hîc omnia persequi, nec lubet, nec vacat. Superest vnum expostulationis nostræ caput, circa nouas nostras quadratrices lineas, quas non ita pridem, vix scilicet ante biennium, inuenimus, nec multò post ad vos misimus: Possem hîc, et sanè potiori jure, eadem verba adjicere quæ vos circa centra grauitatis. Vtinam non misissem. Sed illa nimis acerbam, prorsúsque contumeliosam præ se ferunt exprobrationis speciẽ: quin contra, et misisse lætor, quandoquidem ita vobis placuerunt, et nisi tunc misissem, nunc vtique mitterem. Illas, inquam, lineas ex quibus fiunt spatia plana longitudine infinita, quæ tamen
\note{« æquales sit » : ainsi sur la page. « con\textsuperscript{rum} » : « conicorum » abrégé. « Vtinam non misissem » est écrit en lettres plus grandes, presque capitales : c'est la citation du destinataire.}

\page{70}
\note{Verso du feuillet « 34 », même main ; la phrase de la vue 69 continue. Aucun numéro sur la page. Quelques lettres grises (A, B, C) transparaissent entre les lignes 6--8 depuis l'autre face ou une page voisine ; elles ne sont pas du texte de cette page. Aucune figure n'est tracée sur la page : celle que le texte décrit (« Esto in figura ») est sur la petite bande de la vue 71. Le bas de la page est blanc ; la phrase s'arrête sans ponctuation et continue à la vue 73.}

spatijs finitis vndique clausis sunt æqualia. Vos lineas Roberuallianas, ab inuentoris nomine, vocauistis: ego voco quadratrices, ab earum officio, et inuentionis fine. Ego enim figurarum quadraturæ intentus, dum nihil negligo eorum quæ ad propositum illum finem conducere videntur; præcipuè verò ipsarum figurarum in alias figuras transmutationem experior; in tales lineas incidi, hac ratione.

Esto in figura, trilineam ABC quale requiritur; cujus punctum B sit vertex; recta AB, altitudo; recta AC, basis; et linea BC sit quæcumque curua: nihil enim refert qualiscunque accipiatur. Verùm, vt ex infinitis generibus aliquod hic eligamus, quod vobis instar omnium sit; esto illa curua BC ad easdem partes caua, putà ad partes ductæ rectæ BC; ita vt ipsa tota sit extra triangulum ABC; et eadem à puncto B ad punctum C, continuè recedat à recta BA, et ad rectam CA propiùs accedat; sumpto vtroque, recessu scilicet, et accessu, secundùm perpendiculares à curua BC ad rectas BA, CA, ductas. Tum in ipsa curua BC, sumantur continuè à vertice B, quæcumque et quotcumque puncta D, E, et cæt. à quibus ductæ intelligantur rectæ DF, EG, et cæt. tangentes curuam BC in ijsdem punctis D, E, et cæt. atque occurrentes axi AB producto vltra verticem B, in punctis F, G, et cæt. intelligatur quoque per punctũ C, recta CK tangens eandem curuam BC in puncto C: quæ quidem recta CK vel eidem axi AB occurret vltra verticem B, vel eadem CK eidem AB erit parallela, coincidétque cum recta CR quam ipsi AB ponimus esse parallelam. præterea, à punctis D, E, et cæt. ducantur rectæ DI, EH, axi BA parallelæ, atque occurrentes basi AC in punctis I, H, et cæt. et per punctum A, ipsis tangentibus DF, EG, et cæt. ducantur totidem rectæ ordine parallelæ, AM, quidem ipsi DF; AL autem ipsi EG, et cæt. occurrátque recta AM rectæ DI productæ in M; atque ita habebimus punctum M: occurrat quoque recta AL rectæ EH productæ in L; atque ita rursùs habebimus punctum L. et sic de cæteris. Quo pacto, habebimus à puncto A infinita alia puncta continuo ordine disposita M, L, et cæt. per hæc intelligatur ducta linea continua AML et cæt. illa erit primaria nostra quadratrix: primariam vocamus, quia ipsa prima occurrit, et prima à nobis vulgata est: cæteræ autem ab illa primaria, saltem per occasionem dependerunt. Quòd si tangens CK occurrat axi AB; ductâ rectâ AN parallelâ eidem CK, et productâ rectâ RC donec ipsi AN occurrat in N; erit et punctum N in eadem quadratrice AMLN. Aliàs autem, si CK coincidat cum ipsa CR (cùm scilicet ipsi AB fuerit parallela) linea AML in infinitum producta nunquam concurret cum recta RC etiam infinitè producta; sed hæc RC producta, ipsius AML productæ, erit assymptotos, et punctum N à puncto C infinitè distabit. Potuit etiam loco trilinei, assumi bilineum, aut aliud quodcumque spatium. Sed omnia exequi vnicâ epistolâ, nec possumus, nec volumus; vt ij quibus inuentum placuerit, habeant
\note{« trilineam ABC quale » : ainsi sur la page. Le b de « bilineum » est récrit d'une encre plus noire. Les lettres des points (A, B, C, D, E, F, G, H, I, K, L, M, N, R) sont des capitales posées, plus grandes que le texte.}

\page{71}
\note{Un feuillet blanc de montage, sur lequel est fixée, à droite, une étroite bande de papier portant une figure seule, sans texte. Ses lettres sont des capitales à empattements, d'aspect imprimé ; que la figure soit gravée ou dessinée n'est pas tranché ici. En haut à droite de la bande, d'un trait fin, « 34 bis » (voir l'en-tête). La figure : un long rectangle vertical, côté gauche R Y C N, côté droit K G F B X V A Q P O ; deux verticales intérieures, l'une passant par S, E, H, L, l'autre, plus courte, par T, D, I, M ; les horizontales R \ldots{} (haut), S \ldots, T \ldots, Y B, D X, E V, C H I A, M Q, L P, N O. Une courbe descend de R par S et T jusqu'à B ; une autre monte de C par E et D jusqu'à B ; des droites partent de C vers K et de E vers G, de D vers F ; de A partent les droites vers E et D, et, vers le bas, un faisceau de droites et une courbe passant par M et L jusqu'à N. C'est la figure que décrit le texte de la vue 70 (le trilinéaire ABC, les tangentes DF, EG, CK, les points I, H, M, L, N de la « quadratrix ») ; les lettres S, T, V, X, Y, P, Q, O servent au texte qui suit, à partir de la vue 73.}

\page{73}
\note{Même main. En haut à droite, à l'encre, « 35 » (voir l'en-tête). La vue 72, verso blanc du feuillet de montage de la figure, n'est pas transcrite. Le texte reprend la phrase restée en suspens au bas de la vue 70 (« habeant / quod imitando addere possint »). Le bas de la page est blanc ; la phrase s'arrête sur « et » et continue à la vue 74.}

quod imitando addere possint. Jam ergo, in assumpto exemplo trilinei ABC, positis quæ supra diximus, fit quadrilineum quoddam ABCN duabus curuis BC, AN, et duabus rectis BA, CN, comprehensum; siue id quadrilineum finitum sit versùs N, siue idem in infinitum versùs illam partem abeat: hoc ergo spatium ABCN dico esse trilinei ABC, duplum. Demonstratio nostra omninò vniuersalis erit pro omnibus curuis, et spatijs; poterítque more veterum, per duplicem positionem institui: nos tamen per infinita sic procedemus. Ducantur, aut duci intelligantur à puncto A, ad infinita seu indefinita numero puncta curuæ BC, rectæ AD, AE, et cæt. vt sic spatium ABC in infinita trilinea resolui concipiatur; quæ quidem trilinea totidem rectis AD, AE, et cæt. ac portionibus interceptis curuæ BC comprehendantur. Spatium autẽ ABCN in totidem quadrilinea resoluatur, quot sunt trilinea, quæ quadrilinea à parallelis DM, EL, et cæt. ac portionibus interceptis curuarum BC, AN, constituantur. Erunt ergo singula trilinea cum singulis quadrilineis, super eâdem basi constituta ad puncta D, E, et cæt. propter tangentes, (absque tangentibus enim falsum esset) atque in ijsdem parallelis; putà trilineum ad AD cum quadrilineo ad DM, in ijsdem parallelis DF, MA: trilineam autem ad AE, cum quadrilineo ad EL, in ijsdem parallelis EG, LA, atque ita de reliquis. Quapropter singula quadrilinea singulorum trilineorum erunt vt dupla, ex legibus infiniti; et omnia omnium: hoc est totum spatium ABCN quod ex omnibus quadrilineis constat, duplum erit totius spatij ABC, quod constat ex omnibus trilineis. Patet autem, eodem ratiocinio, quadrilaterum ABDM, trilinei ABD, duplum esse: et quadrilaterum ABEL, trilinei ABE: et sic de cæteris. Si ergo trilineum CAMLN totum extra trilineum ABC existat, vt in assumpto exemplo; erunt duo illa trilinea æqualia; siue punctum N in infinitum abeat; siue non. Quòd si præterèa, eo casu quo curua AMLN tota extra trilineum ABC existit, ex punctis D, E, et cæt. ducantur rectæ DX, EV, basi CA parallelæ atque axi occurrentes in punctis X, V, et cæt. fient spatia BDX, BEV, et cæt. spatijs AIM, AHL, et cæt. singula singulis æqualia. Quoniam enim, ex demonstratione vniuersali præmissa, totum quadrilineum ABDM, totius trilinei ABD, duplum est; et ablatum parallelogrammum AXDI ablati trianguli AXD, est quoque duplum, erit et reliquum reliqui duplum: reliquum autem primum constat ex duobus trilineis BDX, AIM; secundum verò est solum trilineum BDX; quare duo illa trilinea BDX, AIM simul, hujus solius BDX dupla sunt: ac proinde æqualia sunt inter se, trilinea illa BDX, AIM. De cæteris eadem est demonstratio. Sed et trilineum BDF bilineo AM; et
\note{« et » dans « curuis, et spatijs » est ajouté au-dessus de la ligne, de la main du texte. « trilineam » (pour « trilineum ») : ainsi sur la page, comme à la vue 70. Dans « punctum N in infinitum », le N est repassé d'une encre plus noire. Un petit signe gris, en forme de E, entre « existit, » et « ex punctis », est une transparence, non du texte.}

\page{74}
\note{Verso du feuillet « 35 », même main ; la phrase de la vue 73 continue. Aucun numéro sur la page. Le bas de la page est blanc ; la phrase s'arrête sans ponctuation et continue à la vue 75.}

trilineum BEG bilineo AL, æquale esse, facilè demonstrabitur; et multa alia quæ consultò omittimus. Potest quoque ad solida extendi hoc nostrum inuentum; si scilicet, prædictæ omnes figuræ, circa axem AB vtrimque productum quantùm satis, conuertantur: ac spatia quidem solida ad rectas AD, AE, et cæt. constituta, pro pyramidibus; spatia autem solida ad parallelas DM, EL et cæt. pro parallelepipedis accipiantur. Quo pacto, solidum descriptum à quadrilineo ABCN, siue illud versùs N infinitum sit, siue non, triplum erit solidi à trilineo ABC descripti: et solidum à trilineo ACN, in assumpto exemplo, descriptum duplum erit solidi à trilineo ABC descripti. Et hinc habentur innumeræ species solidorum infinitè-finitorum.

Possunt etiam rectæ MI, LH, et cæt. produci versùs puncta D, E, vsque ad puncta T, S, et cæt. ita vt rectæ IT, HS, et cæt. æquales sint rectis DM, EL, et cæt. et per puncta B, T, S, et cæt. potest intelligi curua quadratrix BTS: hæc autem illa erit quam ad vos misimus; de qua ideò, nihil est quod hic addamus: quòd autem illa secundaria sit, manifestum est.

Tandem, ductis tangentibus DF, EG, et cæt. vt supra; potuit, loco puncti A, assumi aliud quodcumque punctum B, vel C, vel quoduis in plano trilinei ABC quantumuis producto, existens \marginal{per quod \uncertain{ducantur} \struck{rectæ} tangentibus illis parallelæ;} quemadmodũ hic ductæ sunt AM, AL, et cæt. et per puncta D, E, et cæt. duci quoque potuerint totidem aliæ rectæ inter se et cuiuis datæ parallelæ, quæ cum tangentibus et tangentium parallelis parallelogramma constituerent, qualia sunt AFDM, AGEL, et cæt. Vnde aliæ infinitæ generabuntur quadratrices. Sed hæc nunc indicasse sufficiat. Vides itaque vir clarissime, quàm latus hoc loco ad imitandũ pateat campus. Vides etiam alia prorsùs à tuis hyperbolicis, diuersa genera solidorum infinitorum, et multitudine innumerabilia, et illis forsan, magis miranda; eo quòd hæc nostra de externa sua latitudine nihil vnquam remittant, vt vestris necessariò accidit. Neque tamen nostra nos ad vestrorum imitationem effinximus (quod si factum fuisset, quantumcumque abstrusa, vobis tamen tribueremus) sed hæc à nostro linearum quadraticum inuento sic dependerunt, vt ab illis sejungi non potuerint. Vides denique nos nec plana, nec solida infinitè-finita præcipuè intendisse; sed nostras quadratrices, quæ ex figurarum in alias transformatione nascuntur; ex quarum origine talia spatia necessariò consecuta sunt; et nobis aliud animo agitantibus, sese vltrò obtulerunt.

Jam, quadratura parabolæ quomodo ex prædictis facilè deducatur, sic ostendimus. Intelligatur in hoc nostro exemplo, curua BC esse quæuis parabola, siue conica illa sit, siue alia: (vnica enim omnibus inseruit demonstratio) cujus axis sit AB; vertex, B; basis AC; et recta BY ipsam tangat in vertice, occurátque rectæ NC
\note{Après « existens », une croix appelle une addition écrite dans la marge gauche, d'une plume plus rapide et plus fine que celle du texte (peut-être une autre main) : « per quod \ldots{} parallelæ; ». Le verbe y est récrit sur lui-même (on lit « duce- », une lettre surchargée, puis « ntur ») ; « rectæ » est barré d'un trait. « quadraticum » : ainsi sur la page (pour « quadratricum »). « occurátque » : ainsi, avec un seul r.}

\page{75}
\note{Même main. En haut à droite de la page, à l'encre, « 36 » (voir l'en-tête) ; au-dessus, sur le bord d'un feuillet qui dépasse la tranche supérieure, « 37 », qui appartient à ce feuillet-là (vue 77) et non à celui-ci. La phrase de la vue 74 continue. Le bas de la page est blanc ; la phrase s'arrête sur une virgule et continue à la vue 76.}

productæ in puncto Y, vt sit parallelogrammum ABYC spatio trilineo parabolico ABC circumscriptum. Ducantur etiam, vel duci intelligantur à singulis punctis curuæ AMLN, putà à punctis M, L, N, et cæt. rectæ MQ, LP, NO, et cæt. basi AC parallelæ occurrentes axi BA producto in punctis Q, P, O, et cæt. quo pacto, constituetur aliud quoddam trilineum ANO, cujus axis erit AO, vertex, A, et basis, NO: in hoc trilineo, rectæ ad axem ordinatim applicatæ erunt MQ, LP, NO, et cæt. quæ ordinatim applicatis in parabola, DX, EV, CA, et cæt. singulæ singulis debito ordine sumptis, erunt æquales; at portiones axis AO inter verticem A, et applicatas interceptæ; putà AQ, AP, AO, et cæt. æquales erunt rectis FX, GV, KA, et cæt. singulæ singulis debito ordine sumptis; quæ omnia ex constructione manifesta sunt. Est autem in quauis parabola, vt FX ad XB, sic GV ad VB, et sic KA ad AB, propter tangentes DF, EG, CK. Quare erit quoque, positâ in nostro exemplo quâuis parabolâ BDEC, vt AQ ad BX, ita AP ad BV, et ita AO ad BA, et cæt. Est ergo curua AMLN parabola ejusdem speciei cum parabola BDEC: cúmque AC, ON, sint æquales; erit spatium AON ad spatium ABC vt axis AO ad axem AB. Ostensum autem est spatium ABC æquale esse spatio ACN; quare spatium AON ad spatium ACN est vt AO ad AB: et componendo, parallelogrammum ACNO ad spatium ACN, siue ad spatiam ABC, se habet vt recta OB ad rectam BA: sed vt parallelogrammum AY ad parallelogrammum AN, ita recta AB ad rectam AO; ergo, ex æquo in ratione perturbata, erit parallelogrammum AY ad spatium ABC vt recta OB ad rectam AO. Datæ autem sunt rectæ illæ OB, AO, quia AO ipsi AK datæ æqualis est, ex constructione: ergo data est ratio parallelogrammi AY ad spatium trilineum parabolicum ABC: vt propositum est: et est talis ratio vt recta composita ex AK, et AB ad rectam AK.

Simili ratiocinio, in solidis ipsarum parabolarum circa axem AB conuersarum, concludemus vniuersaliter sic esse cylindrum AY ad solidum ABC, vt rectam compositam ex AK et dupla ipsius AB ad ipsam eandem AK.

Quomodo ergo in ejusmodi quadratrices inciderim, jam tenes. Quàm verò ingenuè ad vos miserim, ipsi scitis: sciunt et Academiæ nostræ proceres, qui omnes epistolam nostram, antequam ad vos mitteretur, perlegerunt: Sciunt et multi alij cum quibus eandem ego, vel amici communicauimus: sciunt inquam illi omnes me expressis verbis, veluti florem quendam ex horto illo delectum, vobis indicasse quadraturam parabolæ primariæ seu conicæ. Quis igitur meo loco constitutus, fore sperauisset vt clarissimus Toricellius, inde per imitationem, cæteras parabolas quadrandi arreptâ occasione,
\note{Dans « AMLN » (« Est ergo curua AMLN »), le L est récrit par-dessus une autre lettre, d'une encre plus noire. « spatiam ABC » se lit ainsi, pour « spatium ».}

\page{76}
\note{Verso du feuillet « 36 », même main ; la phrase de la vue 75 continue. Aucun numéro sur la page (le « 36 » du recto transparaît, inversé, à gauche). Le bas de la page est blanc ; la phrase s'arrête sur « in » et continue à la vue 77.}

(quod nullius fuit negotij, quia vna eadémque est omnium methodus) hæc verba subjiceret. prædictæ methodi, tum pro quadraturis, tum pro tangentibus, sunt quas minimi præ cæteris ego facio; non tamen patiar mihi illas eripi. |et hæc| Linea Roberualliana, si ortum ducat ex aliqua parabolarum, semper parabola euenit ejusdem speciei, quod ego nouum esse scio licet fortasse turpe videatur hoc fateri. |et rursùs in alia Epistola.| Quadraturas ad clarissimum Roberuallium mitto, fortasse ad subeundam eandem fortunam cum meo centro grauitatis Cycloidis |hoc est Trochoidis. Atque ita, sicuti palam nos accusauerat Toricellius, tamquam si centrum illud nostræ Trochoidis, à nobis illi surreptum fuisset; sic timere se simulauit, ne eodem fato, illæ suæ (si Dijs placet) parabolarum quadraturæ sibi à nobis eriperentur. Quis inquam hoc sperauisset? Nam, Deum immortalem! quid illis in quadraturis aut nouum est, aut ad Toricellium pertinet, vt ei possit eripi? An, in vniuersum, quadraturæ illæ sunt Toricellij? nequaquã: primariæ enim, siue conicæ parabolæ quadratura Archimedis est; cæterarũ autem, Domini de fermat. Dico domini de fermat; quia cæterarum illarum medium à medio Archimedis planè diuersum est, et diuersum esse debuit, quandoquidem ad illas, medium Archimedeum omninò ineptum est. Quòd si omnibus illud aptum fuisset; tunc, quantumuis ab eo diuersum esset medium D. de fermat, omnes tamen illas quadraturas vni Archimedi tribueremus, ac cæteras per imitationem inuentas, ad primariam remitteremus. Siquidem facile est inuentis addere: authorem verò sese præbere, hoc opus, hic labor est. Non igitur aut Toricellij, aut nostræ sunt parabolarum quadraturæ in vniuersum; nec illæ aut ipsi aut nobis eripi possunt. Superest igitur vt de medio decertemus. Sed ad quid hoc? quandò, siue ego vicero, siue Toricellius, ipsa res vel Archimedi cedet, vel D\textsuperscript{o} de fermat. Attamen quod in eo medio præcipuum est, nostrum est, ipso Toricellio concedente, nempe nostra quadratrix, quam ipse Roberuallianam vocat. Quid igitur ipsi relinquitur? forsan, inquiet aliquis, vult Toricellius suum esse, quòd vsus fuerit complementis æqualibus parallelogrammorum, eáque prædictis Roberuallianis quadratricibus accommodauerit, vt duplici positione inscriptorum et circumscriptorum vteretur more veterum. Atqui ob tantillum, quod nec ipsum vniuersale est, adeò sollicitum esse, adeóque inuigilare ne sibi eripiatur, pauperis cujusdam est, qui hoc vnum possideat, non autem ditissimi Toricellij, qui infinitos rerum multò pretiosiorum possidet thesauros. At, dicet alius, Roberuallius vnicam parabolam primariam seu conicam, Toricellius verò omnes omninò quadrauit. Roberuallius scilicet vnicam! Quis autem nos vsque adeò cæcos existimauerit? præcipuè cùm vna eadémque sit omnium methodus, quam suprà ostendimus? Egóne in
\note{Les trois citations des lettres du destinataire (« prædictæ methodi \ldots{} eripi. », « Linea Roberualliana \ldots{} hoc fateri. », « Quadraturas ad clarissimum Roberuallium \ldots{} Cycloidis ») sont écrites d'une écriture plus grande ; entre elles, en petite écriture, les renvois du scripteur, encadrés de traits verticaux, transcrits ici « | … | ». Le trait qui ouvre « |hoc est Trochoidis. » n'a pas de trait fermant. Dans « conicæ » et « parallelogrammorum », une lettre est repassée d'une encre plus noire.}

\page{77}
\note{Même main. En haut à droite, à l'encre, « 37 » (voir l'en-tête ; c'est le chiffre vu au-dessus de la tranche à la vue 75). La phrase de la vue 76 continue. Le bas de la page est blanc ; la phrase s'arrête sur une virgule et continue à la vue 78.}

eo quod difficilius fuit, si tamen quid ibi difficile dici potuit, nempe in quadratricibus ipsis detengendis, atque in primariæ parabolæ quadratura perspicax, in facillimis repentè cæcutiero? Quin ergo saltem enuntiauisti? Satis fuit vnam enuntiare, cæteræ sponte sequebantur. Quid hoc rei est? an tandem ego ea omnia ignorasse censebor, quæcumque vnicâ quam ad Toricellium scripsi epistolâ expressis verbis non comprehendi? Respiciat ille ad verba nostra, vt quid voluerimus intelligat. Florem mittebamus, non arborem. Ac jam decennium est ex quo absolutis nothis illis parabolis, vix animo occurrit, nisi vrgeat occasio, vt illas ampliùs nominem. Toricellio verò ipsæ nouæ sunt, adeóque ipsarum ille non obliuiscitur, vt magnum quid putet, si centum modis illas quadrauerit, cùm tamen infinitis id fieri possit. Rursùs ergo, quid in illis quadraturis nouum est quod ad Toricellium pertineat? non video sanè: attamen scire gestio, ne quod illius est, quódque sibi eripi minimè passurum esse minatur, imprudentes auferamus.

Jam perspiciat quicumque Toricellij legerit epistolas, quàm multa præteream legitimæ expostulationis capita. Enimuerò, illúd ne viro ingenuo ferendum fuit, quod nobis comminando scripsit super aliâ quadam methodo centrorum grauitatis inueniendorum, quam habere se gloriatur? Oro vos, Inquit, ne inter vestra hanc etiã habeatis: nam hoc esset tollere penitùs omne literarum, scientiarúmque commercium. Quid aliud ad manifestum furem scribi potuit? Interim tamen, de illa methodo callidè ac de industria, tacuit Toricellius: ita vt si aliquam ego, aut alius quispiam proferamus, jam ipsi liberum sit illam astutis ejusmodi, atque in longum prospicientibus verbis, sibi asserere, ac de ea locutum esse se, suâ fide affirmare.

Quis rursùs, feret quod ad R. P. Mersennum scribit, cùm de centro nostræ Trochoidis loquitur? Quod certè (ait) imò certissimè scio non habuisse Roberuallium, antequam demonstrationem meam videret; vt P. V. vel ipsemet, vel tandem vniuersa Europa testis esse poterit. De centro illo iam satis supra; imò vsque ad nauseam; nec circa illud vniuersa Europa testis nobis formidanda; quin, si fieri posset, præ cæteris optanda. Verum, quid tale centrum ad vniuersam Europam? Crede mihi clarissime Toricelli: Esto, (quod tamen sine arrogantia dici non potest) quòd in rebus Mathematicis ambo simus egregij; ita vt paucos pares, nullos agnoscamus superiores: nequaquàm tamen, hoc pacto, tales erimus, quos vniuersa respiciat Europa; nempè misellos Geometras de nescio quo puncto disceptantes. Simus potiùs ambo, ego triginta millium peditum nostrorum veteranorum dux, tu totidem vestrorũ: adsit vtrique equitatus tali numero debitus, nihilque desit armorum,
\note{« detengendis » : ainsi sur la page, les lettres « en » repassées d'une encre plus noire. Les deux citations des lettres du destinataire, « Oro vos, Inquit, \ldots{} commercium » et « Quod certè (ait) \ldots{} testis esse poterit », sont d'une écriture plus grande. « P. V. » : « Paternitas Vestra », abrégé ; la citation s'adresse au R. P. Mersenne. « Crede mihi clarissime Toricelli » : le destinataire est nommé ici au vocatif.}

\page{78}
\note{Verso du feuillet « 37 », même main ; la phrase de la vue 77 continue. Aucun numéro sur la page (le « 37 » du recto transparaît, inversé, à gauche). Le bas de la page est blanc ; la phrase s'arrête sans ponctuation et continue à la vue 79.}

annonæ, aut fidei militum erga duces; ac tunc vniuersa forsan nos respiciet Europa.

Hoc loco, vir clarissime, cogitare subit quî fieret vt cùm semel ad te scripserim (prima enim alia nostra de te epistola ad R. P. Mersennum directa fuerat) ídque stilo qui meo et amicorum judicio, nihil omninò acerbi, quamquam post ereptas à te nobis nostras Trochoides, redolet; ipse tamen, e contrario, acri adeò stilo rescripseris; nec mihi soli, quo pacto faciliùs res componerentur; sed tribus (nescio num etiam pluribus) literis ad amplissimos celeberrimósque viros de me scriptis, haud alio argumento quàm quòd existimares (nimis tamen leuiter) centrum Trochoidis ipsius tibi fuisse ereptum. Tantúsne Toricellio earum quas suas putat, nugarum zelus? (liceat eo tibi familiari nugarum vocabulo vti) vt statim atque eas sibi ereptas putauerit,

Irruat et frustra ferro diuerberet vmbras.
\note{Vers isolé, en retrait (cf. Virgile, \emph{Énéide}, VI, 294 ; rapprochement de la transcription).}

Ne quidem cogitando quantas ille, cùm directè, tùm indirectè, ab alijs sumpserit, ob quas periculum sit ne, quamuis placidos, acriùs irritando, ipse vicissim pœnas luat? Atqui consentaneum erat, vir prudens cùm sit, vt meminisset hujus præcepti, quod qui dedit, is procul dubio fuit ad vnguem factus homo; videlicet,

Qui, ne tuberibus proprijs offendat amicum / postulat, ignoscat verrucis illius.
\note{Deux vers isolés, en retrait, sur deux lignes ; la barre oblique marque ici la fin de la première ligne (cf. Horace, \emph{Satires}, I, 3, 73--74 ; rapprochement de la transcription).}

Equidem, inter plurimas hujusce tam acris stili causas, hæc nobis videtur probabilior, quod tu vir clarissime, spatium Mathematicũ ingressus, seu fato, seu sponte, viam à nostris jam ante plures annos tritam inieris, à qua hucvsque parùm deflexeris: vnde non mirum est si in easdem stationes, litora, portus, fluuios, et regiones incidas, quibus illi dudum detectis nomina indiderunt, eáque omnia in chartas intulerunt. Ipse autem, cùm illa à te primùm detecta existimes, fit vt posteà indigneris, si quis contrarium asseruerit, atque id quod verum est candidè enarrauerit. Memineris ergo spatium illud infinitiès infinitè infinitum esse, idémque solidum, imò etiam plusquàm solidum; tibi verò nec pedes, nec pinnas, nec alas deesse: deflectas ergo paululùm vel ad dextram, vel ad sinistram, vel supra, vel infra: curre, nata, vel etiam vola: hæc enim potes omnia, quæ sanè.

pauci, quos æquus amauit / Jupiter, aut ardens euexit ad æthera virtus, / potuere.
\note{Trois lignes isolées, en retrait, la première et la dernière plus courtes (cf. Virgile, \emph{Énéide}, VI, 129--131 ; rapprochement de la transcription). Plus haut, « detectis » : les lettres « te » sont récrites d'une encre plus noire.}

Sic enim fiet, vt, quod non semel, imò pluriès jam præstitisti, et nouas regiones detegas; et viros doctos non solùm adeò feliciter

\page{79}
\note{Même main. En haut à droite, à l'encre, « 38 » (voir l'en-tête). La phrase de la vue 78 continue. Le bas de la page est blanc ; la phrase s'arrête sur « Quòd autem » et continue à la vue 80.}

imiteris, quamquàm nec ipsum laude caret; sed, quod multò laudabilius est, teipsum viris doctis præbeas imitandum.

Hucvsque pro nobis plura diximus: nunc pro diuino Archimede pauca liceat. Bis, vt tua excuses, tantum virum in discrimen adducis vir clarissime. Semel pro libris tuis de motu projectorum; iterùm autem, pro illâ tuâ minimè verâ ratione solidi Trochoidis circa axem, ad suum Cylindrum vt 11 ad 18. Ac primùm quidem, pro libris de motu projectorum hæc ais. Archimedes supposuit olim projecta, non per parabolas, sed per lineas spirales suas procedere. hanc Archimedis suppositionem nullibi videre licuit, in ejus operibus: commentarios autem, forsan, non omnes legi: sed nec eorum authoribus licuit tanto viro absurdas ejusmodi suppositiones affingere. Deinde, pro excusando vestro illo fictitio Trochoidis solido, hæc scribis ad R. P. Mersennum. habemus apud Archimedem prop. 2. de Circuli Dimensione, circulum ad quadratum diametri esse vt 11. ad 14. Quæro ab ipso, (Roberuallio, supple) vndenam putet me habuisse rationem quam ad numeros 11. et 18. reducebam? quæ post verba illa sequitur linea, solitam totius epistolæ redolet acerbitatem. Equidem Archimedes hæc habet: at non dissimulauit statim (nempe prop. 3. quæ manifestò lemma est ad illam secundam) talem rationem 11 ad 14 non esse accuratam, sed tantùm veræ proximam: apud vos autem nihil tale habetur; sed vestram illam rationem 11 ad 18 tanquam accuratam proposuistis, ex inuento priùs centro tanquam accurato deductam. Imò, illam pro accurata exceperunt quicumque existimauerunt vos adeò candidos esse, vt nefas existimaretis ea enuntiare quæ vera non essent. Enimuerò vir clarissime, plerique ex nostris vix persuaderi potuissent Toricellium nobilem adeò Geometram, aliquid purè geometricum, sine demonstratione affirmare voluisse. Sed nec illa vestra ratio 11 ad 18 ex terminis vero proximis ab Archimede assignatis pro circuli dimensione, deducta est; cùm eadem extra ipsos terminos longè euagetur. Vnde non video quid vobis hîc proficiat Archimedis authoritas, præcipuè in materia purè geometrica, vbi pro errore accipitur quidquid accuratè verum non est, quantumcumque illud ad verum proximè accedere deprehendatur.

Hîc fieri posse video, vt aliquis hujusce nostræ epistolæ stilum ideò carpat, quòd ille, nec amico, nec aduersario conuenire videatur; vt potè qui pro amico, acrior, pro aduersario contra, lenior quàm par sit appareat. Equidem, clarissimum Toricellium aduersarium habere, absit vt vnquam optauerim: aduersarius sanè illi ego ero nunquam, nisi ipse prior talem me effecerit. Quòd autem
\note{Les deux citations des lettres du destinataire, « Archimedes supposuit \ldots{} procedere. » et « habemus apud Archimedem \ldots{} reducebam? », sont d'une écriture plus grande ; « (Roberuallio, supple) », écrit de la même grande écriture, est la glose du scripteur. « vndenam » est écrit en un mot.}

\page{80}
\note{Verso du feuillet « 38 », même main ; la phrase de la vue 79 continue. Aucun numéro sur la page (le « 38 » du recto transparaît, inversé, à gauche). Le bas de la page est blanc sous la dernière ligne.}

amicum et cupierim et adhuc cupiam, argumentum certissimum est, quòd prior amauerim, ac nomen ejus celebre, per Galliam, quàm maximè potui, reddiderim. Siccine ergo (vrgebit censor) cum amicis tuis te gerere solitus es? Primùm quidem, apologiam contra acerbam ipsius accusationem mihi debui. Deinde metui (fateor) ne ipse, quem summoperè amicum mihi cupio, ex illis esset qui aliena veluti perspicillis cauis respiciunt; sua, conuexis, aut ijs forsan, quæ plurimis faciebus distinguntur: vnde fit vt ijdem aliena contractiora; sua verò ampliora, aut numerosiora, aut etiam pulchris coloribus ornatiora quam sint reuerà videre videantur. Itaque admonere eum volui charitatis, vt amorem proprium alieno temperaret. Ac, ne ad excitandum duriusculus haberetur, stilum adhibui vtcumque acutum et mordacem. Sic enim fore speraui vt, sapiens cùm sit, se ab amante pungi sentiret, atque ita ad redamandum acriùs incitaretur. Quamquam autem tot paginas minimè inutiles fore spero; doleo tamen quòd illas in tractando ejusmodi ingrato, ac planè tædioso argumento insumere oportuerit; cùm alia ferè innumera longè suauiora, ac viris doctis, vt puto, acceptiora, cùm ex nobis, tùm ex nostris habeamus; qualia sunt quæ sequuntur.

Circa Analysim quidem. De æquationum recognitione, et emendatione, nouâ prorsùs methodo; de earundem determinatione; ac de ipsarum per locos proprios resolutione, atque compositione. Circa Geometriam. de Locis planis, solidis, atque ad superficiem; vbi in specie, restituta habemus loca solida ad tres et quatuor lineas: de Cylindris, et conis isoperimetris, cùm demptâ base, tum additâ: de ijsdem sphæræ inscriptis, et circumscriptis: seu spatiorum solidorum, seu etiam superficierum tantùm habeatur ratio: vbi miraberis forsan quâ ratione à nobis concludi potuerit, positâ sphæræ diametro 32 partiũ, axem coni inscripti cujus superficies, comprehensâ base, sit maxima, esse hanc apotomen $23 - \sqrt{17}$. Si Sphæræ superficies vno, duobúsue, vel tribus, aut pluribus circulis, in quotcunque et quascumque portiones secta sit; quamcumque ex illis portionibus cum alia, ac cum tota comparamus; ac vniuscujusque centrum grauitatis assignamus. Circa cylindricas, et conicas superficies scalenas, tum etiam circa rectas, mira habemus; inter illa perpende qualenam sit hoc problema. Portionem superficiei cylindri recti exhibemus, quæ superficiei datæ cylindri scaleni sit æqualis: sed et istud. Dato quadrato, æqualem damus cylindricæ superficiei portionem, idque absolutè, nullâ suppositâ circuli quadraturâ, et exclusis cylindri basibus. Problemata, atque theoremata innumera habemus soluta, cùm circa conicas sectiones, tùm circa alia ferè omnia geometriæ hucvsque nota tam theoretica, quã practicæ capita. Circa Arithmeticam, Musicam, Opticam, Astronomiam
\note{« $\sqrt{17}$ » : le signe de la racine est un petit trait en forme de r, placé devant 17 ; il est rendu par $\sqrt{\ }$. Dans la marge gauche, à la hauteur de « axem coni inscripti », un long trait horizontal, d'une encre plus pâle (marque de lecteur ?). « Astrono- » est coupé en fin de ligne ; « miam » est écrit seul sous la ligne, contre la marge droite, au bas du texte : c'est la fin du mot, et sans doute la réclame de la page suivante, hors de ce lot. La phrase continue au-delà de la vue 80.}

\end{document}
